\chapter{HOS coupling and currents}
\label{ch:hos}

\section{Principle}
High-Order Spectral (HOS) solvers compute nonlinear, phase-resolved wave fields over large
domains and long durations at a fraction of the cost of CFD. HOS-NWT models a wave tank with
a wavemaker~\cite{ducrozet2012modified}; HOS-Ocean models a periodic open-sea
domain~\cite{ducrozet2016hosocean}. Both write their solution as spectral modes in the file
\file{modes\_HOS\_SWENSE.dat}.

The Grid2Grid library~\cite{choi2019grid2grid} reads this file and reconstructs the free-surface
elevation $\eta$, the velocity $\mathbf u$ and the dynamic pressure $p_d$ at any point and time.
In \fwt{}, Grid2Grid is used as an \emph{external wave forcing} of \wtwof{}: the relaxation
zones blend the CFD solution towards the HOS solution, exactly as they do for the algebraic wave
theories (\wtwof{} manual, Sec.~2.7). The coupling strategy and its validation are described
in~\cite{aliyar2022numerical}.

\section{Workflow}
\begin{enumerate}
  \item Run HOS-NWT or HOS-Ocean for the wave condition. Keep
        \file{modes\_HOS\_SWENSE.dat}.
  \item Copy \file{modes\_HOS\_SWENSE.dat} into the case folder \file{constant/}.
  \item Describe the HOS file in \file{constant/Grid2Grid.dict} (Sec.~\ref{sec:g2gdict}).
  \item Select the HOS forcing in \file{constant/waveProperties.input}
        (Sec.~\ref{sec:hos-setup}).
  \item Load the libraries in \file{system/controlDict}:
\begin{lstlisting}
libs ( "libwaves2Foam.so" "libGrid2Grid.so" );
\end{lstlisting}
  \item Prepare and run the case with the usual \wtwof{} sequence:
\begin{lstlisting}[style=shell,caption={Running an HOS-coupled case (shortened from \file{tutorials/waveFoam/Allrun}).}]
blockMesh
waveGaugesNProbes
setWaveParameters
relaxationZoneLayout
setWaveField          # initial field from HOS at the start time
waveFoam              # or waveIsoFoam
\end{lstlisting}
\end{enumerate}

\NB{\app{setWaveField} initialises $\alpha$ and $\mathbf u$ from the HOS solution at the start
time of the simulation (\kw{startTime} plus \kw{referenceTime}). It advances Grid2Grid to that
time first, so a simulation can start at any time inside the HOS record.}

\section{The Grid2Grid dictionary}
\label{sec:g2gdict}
\begin{lstlisting}[caption={\file{constant/Grid2Grid.dict}.},label={lst:g2gdict}]
Grid2Grid       HOSNWT;           // name of the sub-dictionary below

HOSNWT
{
    type                NWT;      // NWT (HOS-NWT) or Ocean (HOS-Ocean)
    fileType            ASCII;    // ASCII or HDF5
    filePath            constant/modes_HOS_SWENSE.dat;
    interpolationOrder  3;        // 3: cubic spline (default)
    zMeshType           meshRatio;// uniform, sine or meshRatio
    zMin                -5;       // lowest z of the reconstruction grid [m]
    zMax                20;       // highest z (above the highest crest)
    nZMin               50;       // points below z = 0
    nZMax               50;       // points above z = 0
    zMinRatio           3;        // grid stretching below z = 0 (meshRatio)
    zMaxRatio           3;        // grid stretching above z = 0 (meshRatio)
}
\end{lstlisting}
Grid2Grid reconstructs the HOS field on a vertical grid between \kw{zMin} and \kw{zMax} and
interpolates from it. \kw{zMin} must reach at least the bottom of the CFD domain and \kw{zMax}
the top of the highest crest. With \kw{meshRatio}, the points are clustered towards $z=0$.
The full list of options is in the Grid2Grid manual~\cite{choi2019grid2grid}.

\section{Selecting the HOS forcing}
\label{sec:hos-setup}
\begin{lstlisting}[caption={HOS forcing in \file{constant/waveProperties.input}.},label={lst:hosprops}]
seaLevel            0.0;
relaxationNames     (inlet outlet);
initializationName  inlet;

externalForcing     Grid2Grid;

Grid2GridCoeffs
{
    waveType                Grid2Grid;
    Grid2GridInputFileName  Grid2Grid.dict;
    referencePosition       (-16.95 -14.84 0);
    referenceTheta          0;      // degrees
    referenceTime           0;      // seconds
    Tsoft                   0;      // seconds
}

inletCoeffs
{
    waveType        externalSource; // take the wave from externalForcing
    relaxationZone
    {
        relaxationScheme Spatial;
        relaxationShape  Rectangular;
        relaxType        INLET;
        startX           (-10.11 0 -5);
        endX             (-2 0 2);
        orientation      (1 0 0);
    }
}
// outletCoeffs: the same, with relaxType OUTLET
\end{lstlisting}

The keyword \kw{externalForcing} selects the HOS wave field for the whole case; every
relaxation zone that should follow it uses \kw{waveType externalSource}. A relaxation zone may
also use any other \wtwof{} theory. The relaxation zone settings are those of \wtwof{}
(\wtwof{} manual, Sec.~2.7.2).

\begin{table}[htbp]
\centering\small
\begin{tabularx}{\linewidth}{@{}l l X@{}}
\toprule
Keyword & Default & Meaning \\
\midrule
\kw{Grid2GridInputFileName} & -- (required) & Grid2Grid dictionary in \file{constant/} \\
\kw{referencePosition} & $(0\;0\;0)$ & OpenFOAM point that coincides with the HOS origin \\
\kw{referenceTheta} & 0 & rotation of the HOS axes about $z$, in degrees \\
\kw{referenceTime} & 0 & HOS time at OpenFOAM time $t=0$ \\
\kw{Tsoft} & 0 & ramp-up time of the HOS field; 0 = no ramp \\
\bottomrule
\end{tabularx}
\caption{Keywords of \kw{Grid2GridCoeffs}.}
\label{tab:g2gcoeffs}
\end{table}

\Warn{Only the keywords in Table~\ref{tab:g2gcoeffs} are read. Some older cases contain
\kw{refTime} or \kw{refDirection}; these are ignored. Use \kw{referenceTime} to shift the HOS
time, and \kw{referenceTheta} to rotate the HOS field.}

\subsection{Mapping between OpenFOAM and HOS}
A point $\mathbf x$ and time $t$ in OpenFOAM are mapped to the HOS coordinates as
\begin{gather}
  \begin{pmatrix} x_H \\ y_H \end{pmatrix} =
  \begin{pmatrix} \cos\theta & -\sin\theta \\ \sin\theta & \cos\theta \end{pmatrix}
  \begin{pmatrix} x - x_\mathrm{ref} \\ y - y_\mathrm{ref} \end{pmatrix},
  \\
  z_H = z - z_\mathrm{ref} - z_\mathrm{sea},
  \qquad
  t_H = t + t_\mathrm{ref},
\end{gather}
with $\theta = \kw{referenceTheta}$, $\mathbf x_\mathrm{ref}=\kw{referencePosition}$,
$t_\mathrm{ref}=\kw{referenceTime}$ and $z_\mathrm{sea}=\kw{seaLevel}$. The
returned velocity is rotated back to the OpenFOAM axes, \kw{seaLevel} is added to $\eta$, and
the hydrostatic pressure at \kw{seaLevel} is added to $p_d$.

Gravity in \file{constant/g} must point in the negative $z$ direction; otherwise the solver
stops with an error.

\subsection{Ramp-up}
With $\kw{Tsoft}>0$, $\eta - \texttt{seaLevel}$, $\mathbf u$ and $p_d$ from HOS are multiplied by
\begin{equation}
  f(t) = \sin\!\left(\frac{\pi}{2}\,\frac{\min(t,\,T_\mathrm{soft})}{T_\mathrm{soft}}\right),
\end{equation}
which rises from 0 to 1 over the first $T_\mathrm{soft}$ seconds. A ramp is useful when the
simulation starts from still water; when the initial field is set by \app{setWaveField}, use
$\kw{Tsoft}=0$.

\section{HOS waves with a current}
\label{sec:hos-current}
\kw{Grid2GridWithCurrent} adds a steady current $\mathbf U_c(z)$ to the HOS velocity. It reads
all keywords of Table~\ref{tab:g2gcoeffs} from the same \kw{Grid2GridCoeffs} dictionary, plus
those in Table~\ref{tab:current}.
\begin{lstlisting}[caption={HOS waves with a depth-varying current (\file{tutorials/HOS\_with\_Current}).}]
externalForcing     Grid2GridWithCurrent;

Grid2GridCoeffs
{
    waveType                Grid2GridWithCurrent;
    Grid2GridInputFileName  Grid2Grid.dict;
    referencePosition       (-16.95 -14.84 0);
    referenceTheta          0;
    referenceTime           0;
    Tsoft                   0;

    profileType     depthVarying;   // constant | depthVarying | parabolic
    Uc              (0 0 0);        // current below ZC
    UF              (0.25 0 0);     // current at the free surface
    ZC              -5;             // transition depth
    currentTsoft    10;             // ramp-up time of the current
}
\end{lstlisting}
The HOS wave field itself is \emph{not} modified by the current; the current is superposed on the
HOS velocity. This is a valid model when the HOS simulation was run without current and the
current is weak compared with the wave phase speed.

\subsection{Current profiles}
\label{sec:current-profiles}
\begin{description}
  \item[\kw{constant}] $\mathbf U_c = \kw{Uc}$ everywhere.
  \item[\kw{depthVarying}] $\mathbf U_c = \kw{Uc}$ below $z = \kw{ZC}$, and a linear transition
    to \kw{UF} at the still-water level:
    \begin{equation}
      \mathbf U_c(z) = (1-w)\,\mathbf U_\mathrm{c} + w\,\mathbf U_F,\qquad
      w = \frac{z - Z_C}{\texttt{seaLevel} - Z_C},\quad z > Z_C .
    \end{equation}
  \item[\kw{parabolic}] a profile along $x$ with its maximum $|\mathbf U_F|$ at $z=z_0$,
    falling to zero at $z_0 \pm H$:
    \begin{equation}
      U_{c,x}(z) = \max\!\left(|\mathbf U_F|\left[1 - \left(\frac{z-z_0}{H}\right)^2\right],\,0\right).
    \end{equation}
\end{description}
The current is ramped up with
$\sin\left(\frac{\pi}{2}\min(t,T_c)/T_c\right)$, $T_c=\kw{currentTsoft}$.

\begin{table}[htbp]
\centering\small
\begin{tabularx}{\linewidth}{@{}l l X@{}}
\toprule
Keyword & Default & Meaning \\
\midrule
\kw{profileType}  & \kw{constant} & \kw{constant}, \kw{depthVarying} or \kw{parabolic} \\
\kw{Uc}           & $(0\;0\;0)$ & current below \kw{ZC} (and everywhere for \kw{constant}) \\
\kw{UF}           & $(0\;0\;0)$ & current at the still-water level; peak value for \kw{parabolic} \\
\kw{ZC}           & 0 & transition depth for \kw{depthVarying} \\
\kw{currentTsoft} & 0 & ramp-up time of the current; 0 = no ramp \\
\kw{H}            & $|\kw{ZC}|$ & half-width of the parabolic profile \\
\kw{z0}           & $(\texttt{seaLevel}+\kw{ZC})/2$ & height of the parabolic maximum \\
\bottomrule
\end{tabularx}
\caption{Current keywords of \kw{Grid2GridWithCurrent}.}
\label{tab:current}
\end{table}

\section{Current without waves}
\label{sec:current-only}
The wave theory \kw{DepthvaryingCurrent} imposes only a current, with the same three profiles.
It is used as the \kw{waveType} of a relaxation zone and is set up like any \wtwof{} theory:
\begin{lstlisting}[caption={Current-only relaxation zone (\file{tutorials/CurrentWithProfile}).}]
inletCoeffs
{
    waveType        DepthvaryingCurrent;
    profileType     depthVarying;   // constant | depthVarying | parabolic
    U               (0 0 0);        // current below ZC
    UF              (1.0 0 0);      // current at the free surface
    ZC              -5.0;
    Tsoft           5;              // ramp-up time
    localSeaLevel   0.0;            // default: seaLevel
    relaxationZone { ... }
}
\end{lstlisting}
\NB{In \kw{DepthvaryingCurrent} the current below \kw{ZC} is called \kw{U}, and
\kw{profileType}, \kw{U}, \kw{UF}, \kw{ZC} and \kw{Tsoft} are required. In
\kw{Grid2GridWithCurrent} the same quantity is called \kw{Uc}, and all current keywords have
defaults. The profiles and the ramp are otherwise identical (Sec.~\ref{sec:current-profiles}),
with \kw{localSeaLevel} in place of \kw{seaLevel}.}
